\documentclass[a4paper, serif, 10pt]{moderncv}

%% moderncv
% color and style
\moderncvcolor{black}
\moderncvstyle{banking}

% renew moderncv command to adapt for CJK colon
\renewcommand*{\cvitem}[3][.25em]{%
  \ifstrempty{#2}{}{\hintstyle{#2}：}{#3}%
  \par\addvspace{#1}}

\renewcommand*{\cvitemwithcomment}[4][.25em]{%
  \savebox{\cvitemwithcommentmainbox}{\ifstrempty{#2}{}{\hintstyle{#2}：}#3}%
  \setlength{\cvitemwithcommentmainlength}{\widthof{\usebox{\cvitemwithcommentmainbox}}}%
  \setlength{\cvitemwithcommentcommentlength}{\maincolumnwidth-\separatorcolumnwidth-\cvitemwithcommentmainlength}%
  \begin{minipage}[t]{\cvitemwithcommentmainlength}\usebox{\cvitemwithcommentmainbox}\end{minipage}%
  \hfill% fill of \separatorcolumnwidth
  \begin{minipage}[t]{\cvitemwithcommentcommentlength}\raggedleft\small\itshape#4\end{minipage}%
  \par\addvspace{#1}}

%% page layout/margins
\usepackage[top=2.5cm, bottom=2.5cm, left=1.5cm, right=1.5cm]{geometry}

\nopagenumbers{}

%% Babel config for English language
\usepackage[english]{babel}

%% fontspec
\usepackage{fontspec}

\IfFontExistsTF{Linux Libertine}{
  \setmainfont[Ligatures={TeX, Common}, Numbers=Lining, ItalicFont=Linux Libertine]{Linux Libertine}
}{}
\IfFontExistsTF{Linux Libertine O}{
  \setmainfont[Ligatures={TeX, Common}, Numbers=Lining, ItalicFont=Linux Libertine O]{Linux Libertine O}
}{}

%% CTeX
% CJK support, used to show CJK characters in the resume
%
% - heading=false: disable ctex heading
% - punct=kaiming: use kaiming punctuations styles for CJK
% - scheme=plain: use plain scheme, do not override \normalsize font size
% - space=auto: space settings for CJK characters
%
% CJK fonts are resolved per locale language, e.g. Noto Serif/Sans CJK JP
% for Japanese resumes, with fallback to the next region if not installed
%
% ref:
% - http://ctan.mirrorcatalogs.com/language/chinese/ctex/ctex.pdf
\usepackage[UTF8, heading=false, punct=kaiming, scheme=plain, space=auto]{ctex}

\IfFontExistsTF{Noto Serif CJK TC}{
  \setCJKmainfont{Noto Serif CJK TC}
}{}
\IfFontExistsTF{Noto Sans CJK TC}{
  \setCJKsansfont{Noto Sans CJK TC}
}{}

%% line spacing
\usepackage{setspace}
\setstretch{1.125}

%% URL styling - use normal text instead of monospace
\urlstyle{same}

%% Patch \httplink and \httpslink to handle full URLs
% The original moderncv macros always prepend http:// or https:// to URLs.
% This causes issues when the URL already has a protocol (e.g., https://example.com)
% resulting in malformed URLs like https://https://example.com.
% This patch checks if the URL already contains "://" and uses it directly if so.
% Note: We use \str_set:Nx (x-expansion) to expand macros like \@homepage first.
\ExplSyntaxOn
\str_new:N \g_yamlresume_url_str

\RenewDocumentCommand{\httpslink}{O{}m}{
  \str_set:Nx \g_yamlresume_url_str {#2}
  \str_if_in:NnTF \g_yamlresume_url_str {://}
    { \url{#2} }
    { \url{https://#2} }
}

\RenewDocumentCommand{\httplink}{O{}m}{
  \str_set:Nx \g_yamlresume_url_str {#2}
  \str_if_in:NnTF \g_yamlresume_url_str {://}
    { \url{#2} }
    { \url{http://#2} }
}
\ExplSyntaxOff

\name{林雅婷}{}
\title{人力資源經理｜人才策略與組織發展}
\phone[mobile]{+886 912 345 678}
\email{yating.lin@example.com}
\homepage{https://www.linkedin.com/in/yating-lin-hr}

\address{中國臺灣，臺灣，臺北市，臺北市信義區松高路 88 號 12 樓，110}{}{}

\extrainfo{{\small \faLinkedin}\ \href{https://www.linkedin.com/in/yating-lin-hr}{@yating-lin-hr} {} {} {} • {} {} {} 
{\small \faMedium}\ \href{https://medium.com/@yatinglinhr}{@yatinglinhr}}

\begin{document}

\maketitle

\section{簡介}

\cvline{}{資深人力資源經理，擁有 10 年以上科技與顧問產業經驗，專長於人才策略、組織發展、績效管理與員工關係。
%
\begin{itemize}
\item 曾領導 8 人人力資源團隊，服務超過 1,200 名員工，成功將年度離職率降低 18\%。
\item 主導人才盤點與接班計畫，建立關鍵職位發展藍圖，提升內部晉升比例至 65\%。
\item 熟悉勞動法令與薪酬架構設計，能平衡員工體驗與組織營運目標。
\item 具跨國顧問與校園招募經驗，擅長與高階管理層溝通並推動變革。
\end{itemize}}
\section{教育背景}

\cventry{2012年9月–2014年6月}
        {碩士，人力資源管理研究所，成績：4.1/4.3}
        {國立臺灣大學}
        {\url{https://www.ntu.edu.tw}}
        {}
        {\begin{itemize}
\item 主修人力資源管理，專注於策略性人力資源、組織行為與薪酬制度設計。
\item 完成碩士論文《高科技產業人才留任因素之研究》，以問卷調查分析 300 位工程師樣本。
\item 擔任人力資源管理研究所學會幹部，籌辦企業參訪與專題講座。
\end{itemize}
\textbf{課程}：策略性人力資源管理、組織行為、薪酬管理、績效管理、訓練與發展、勞動法令、研究方法、人力資源分析}

\cventry{2008年9月–2012年6月}
        {學士，企業管理學系，成績：3.8/4.0}
        {國立政治大學}
        {\url{https://www.nccu.edu.tw}}
        {}
        {\begin{itemize}
\item 建立企業管理基礎，涵蓋行銷、財務、作業管理與人力資源等領域。
\item 參與企業實習與個案競賽，培養商業分析與團隊協作能力。
\end{itemize}
\textbf{課程}：管理學、經濟學、會計學、行銷管理、財務管理、作業管理、人力資源管理、商業分析}
\section{工作經歷}

\cventry{2021年3月至今}
        {人力資源經理}
        {雲暉科技股份有限公司}
        {\url{https://www.yunhui-tech.com}}
        {}
        {\begin{itemize}
\item 負責 1,200 人規模科技公司之全方位人力資源策略，涵蓋招募、薪酬、績效、培訓與員工關係。
\item 領導 8 人人力資源團隊，建立部門 OKR 與月度營運檢視機制，提升專案交付效率 25\%。
\item 重新設計職等職級與薪酬架構，完成市場薪資對標，關鍵人才留任率提升至 92\%。
\item 推動年度人才盤點與接班梯隊計畫，辨識 45 位高潛力人才並規劃個人發展藍圖。
\item 建立雇主品牌與校園招募策略，年度招募週期縮短 20\%，並提升新人到職率 15\%。
\item 處理勞資爭議與勞動檢查，確保公司遵循勞基法及相關法令，達成零重大裁罰紀錄。
\end{itemize}
\textbf{關鍵字}：人才策略、組織發展、績效管理、薪酬設計、員工關係、招募甄選、勞動法令、HRIS}

\cventry{2016年7月–2021年2月}
        {資深人力資源顧問}
        {智聯國際顧問股份有限公司}
        {\url{https://www.zhilian-consulting.com}}
        {}
        {\begin{itemize}
\item 提供科技、製造與零售客戶組織設計、職等職級與薪酬制度顧問服務，累積 30 個以上專案。
\item 主導人才盤點與接班計畫，協助客戶建立關鍵職位勝任力模型與發展路徑。
\item 設計訓練發展藍圖與講師培育制度，協助客戶將內部訓練時數提升 40\%。
\item 協助企業因應勞動法令修法，進行風險盤點、制度修訂與主管培訓。
\item 管理 3 人專案團隊，負責提案簡報、專案時程、預算控制與客戶關係維護。
\end{itemize}
\textbf{關鍵字}：人力資源顧問、組織設計、職等職級、人才盤點、訓練發展、勞動法令、專案管理}

\cventry{2014年8月–2016年6月}
        {人力資源專員}
        {宏亞零售集團}
        {\url{https://www.hongya-retail.com}}
        {}
        {\begin{itemize}
\item 負責零售門市人員招募、考勤薪酬、教育訓練與員工活動規劃。
\item 建立招募漏斗與面試評分表，將門市人員平均招募天數由 35 天降至 24 天。
\item 執行新人訓練與在職訓練，協助 200 位以上員工完成年度必修課程。
\item 規劃員工關懷與家庭日活動，提升員工滿意度 12\%。
\end{itemize}
\textbf{關鍵字}：招募、考勤薪酬、教育訓練、員工活動、零售業}
\section{語言}

\cvline{中文}{母語或雙語水平 \hfill \textbf{關鍵字}：母語、商務書寫}
\cvline{英語}{完全專業水平 \hfill \textbf{關鍵字}：TOEIC 950、IELTS 7.5、跨國會議}
\cvline{日語}{初級水平 \hfill \textbf{關鍵字}：JLPT N4}
\section{專業技能}

\cvline{人力資源策略}{專家 \hfill \textbf{關鍵字}：組織設計、人才盤點、接班計畫、人力預算、變革管理}
\cvline{招募與雇主品牌}{高級 \hfill \textbf{關鍵字}：校園招募、面試技巧、雇主品牌、招募漏斗、ATS}
\cvline{薪酬與績效管理}{高級 \hfill \textbf{關鍵字}：薪酬架構、職等職級、KPI、OKR、獎酬設計}
\cvline{勞動法令與員工關係}{高級 \hfill \textbf{關鍵字}：勞基法、勞資爭議、員工關係、合規管理}
\cvline{人力資源資訊系統}{中級 \hfill \textbf{關鍵字}：SAP SuccessFactors、Workday、Power BI、Excel}
\cvline{教育訓練與發展}{中級 \hfill \textbf{關鍵字}：TNA、訓練藍圖、講師培育、學習平台}
\section{榮譽}

\cventry{2023年11月}
        {中華人力資源管理協會}
        {年度最佳人力資源團隊獎}
        {}
        {}
        {\begin{itemize}
\item 表揚在人才策略、員工體驗與組織變革具卓越成效之人資團隊。
\item 帶領團隊完成薪酬架構改革與雇主品牌計畫，獲得評審一致肯定。
\end{itemize}}

\cventry{2022年9月}
        {亞太人力資源評論}
        {卓越雇主品牌獎}
        {}
        {}
        {\begin{itemize}
\item 肯定公司在校園招募、員工價值主張與多元共融政策之表現。
\item 成功提升品牌知名度與新人到職率，並降低招募成本。
\end{itemize}}
\section{證書}

\cventry{2018年3月}
        {中華人力資源管理協會}
        {全方位人力資源管理師}
        {\url{https://www.chrma.org.tw/certification/hr-manager}}
        {}
        {}

\cventry{2019年6月}
        {國際專案管理學會}
        {國際專案管理師 PMP}
        {\url{https://www.pmi.org/certifications/project-management-pmp}}
        {}
        {}

\cventry{2020年10月}
        {台灣勞動法學會}
        {勞動法令遵循專家}
        {\url{https://www.taiwanlaborlaw.org/certification/compliance}}
        {}
        {}
\section{出版刊物}

\cventry{2022年12月}
        {數位轉型下的人才策略：以台灣科技業為例}
        {人力資源管理學報}
        {\url{https://doi.org/10.1234/hrmj.2022.001}}
        {}
        {\begin{itemize}
\item 探討數位轉型對科技業人才策略、組織設計與關鍵職能之影響。
\item 以三家台灣上市科技公司為個案，提出人才培育與接班規劃建議。
\end{itemize}}

\cventry{2021年5月}
        {從員工體驗出發的留才策略}
        {亞太人力資源評論}
        {\url{https://www.aphrreview.com/articles/employee-experience-retention}}
        {}
        {\begin{itemize}
\item 分析員工體驗對留任意願之影響，提出五大員工旅程改善方向。
\item 提供人資主管可執行之留才措施與衡量指標。
\end{itemize}}
\section{項目}

\cventry{2022年1月–2023年12月}
        {為關鍵職位建立人才盤點、發展藍圖與接班梯隊，確保組織中長期人才供給。}
        {人才發展與接班計畫}
        {\url{https://www.yunhui-tech.com/esg/talent-development}}
        {}
        {\begin{itemize}
\item 盤點 12 個關鍵職位與 45 位高潛力人才，建立九宮格人才地圖。
\item 設計個人發展計畫與導師制度，兩年內內部晉升比例提升至 65\%。
\item 建立接班梯隊檢視機制，每季向高階管理層報告人才準備度。
\end{itemize}
\textbf{關鍵字}：人才盤點、接班梯隊、發展藍圖、關鍵職位}

\cventry{2021年3月–2021年11月}
        {導入一站式員工服務平台，整合入離職、考勤、薪酬查詢與意見回饋流程。}
        {員工體驗平台導入}
        {\url{https://www.yunhui-tech.com/hr/employee-experience}}
        {}
        {\begin{itemize}
\item 與 IT 及外部供應商協作，完成需求訪談、流程設計與系統上線。
\item 將常見人資問題處理時間縮短 40\%，員工服務滿意度提升至 88\%。
\item 建立人資數據儀表板，追蹤服務量、處理時效與員工回饋趨勢。
\end{itemize}
\textbf{關鍵字}：員工體驗、流程自動化、HRIS、服務中心}
\section{興趣愛好}

\cvline{人力資源社群}{專業分享、講座、讀書會}
\cvline{運動}{瑜伽、慢跑、登山}
\cvline{閱讀}{管理學、心理學、人物傳記}
\section{志願服務}

\cventry{2019年1月至今}
        {志工講師}
        {中華人力資源管理協會}
        {\url{https://www.chrma.org.tw}}
        {}
        {\begin{itemize}
\item 擔任人力資源管理師認證課程講師，分享薪酬設計與員工關係實務。
\item 協助協會辦理年度論壇與企業參訪，促進人資專業交流。
\end{itemize}}
    

\cventry{2020年3月–2023年12月}
        {職涯導師}
        {台灣青年職涯發展協會}
        {\url{https://www.twyouthcareer.org.tw}}
        {}
        {\begin{itemize}
\item 提供青年一對一職涯諮詢，協助履歷健檢、面試準備與職涯規劃。
\item 參與校園職涯講座，分享人資產業趨勢與求職技巧。
\end{itemize}}
    
\end{document}